\chapter{근을 담는 새로운 수의 세계}
\label{chap:v2-motivation}

\begin{learninggoal}
이 장을 마치면 다음을 설명할 수 있다.
\begin{itemize}
  \item 왜 하나의 체 안에서 모든 다항방정식을 풀 수 없는가?
  \item 근을 첨가한다는 말을 몫체로 어떻게 구현하는가?
  \item 체확장의 차수와 자동동형이 각각 무엇을 측정하는가?
  \item 체확장론이 갈루아 이론으로 어떻게 이어지는가?
\end{itemize}
\end{learninggoal}

\section{방정식은 계수체를 벗어나기도 한다}

다항식의 계수가 어떤 체 $K$에 속한다고 해서 그 근도 반드시 $K$에 속하는 것은 아니다. 예를 들어
\[
  X^2-2\in\Q[X]
\]
는 유리수 근을 갖지 않고,
\[
  X^2+1\in\R[X]
\]
는 실수 근을 갖지 않는다. 또한
\[
  X^2+X+1\in\F_2[X]
\]
도 $0$과 $1$ 어느 쪽도 근으로 갖지 않는다.

이때 선택지는 두 가지다. 첫째, 방정식을 포기한다. 둘째, 근을 담을 수 있도록 수의 세계를 넓힌다. 대수학은 두 번째 길을 택한다.

\begin{keyidea}
체확장은 기존의 계산법을 버리는 과정이 아니다. 원래 체의 덧셈과 곱셈을 그대로 보존하면서, 필요한 원소만 새로 첨가하는 과정이다.
\end{keyidea}

\section{근을 첨가하는 가장 작은 체}

$X^2-2$의 한 근을 $\alpha$라고 부르고 $\alpha^2=2$라는 관계를 강제하자. 그러면 사칙연산으로 만들어지는 원소는
\[
  a+b\alpha\qquad(a,b\in\Q)
\]
의 꼴로 정리된다. 이 집합은 실제로 체이며, 보통 $\Q(\alpha)$ 또는 $\Q(\sqrt2)$로 쓴다.

이 구성을 다항식환의 언어로 쓰면 더 깔끔하다. $X^2-2$는 $\Q[X]$에서 기약이므로
\[
  \Q[X]/(X^2-2)
\]
는 체이다. 여기서 $\overline X=X+(X^2-2)$라 두면
\[
  \overline X^2=2
\]
가 성립한다. 즉 변수 $X$의 잉여류가 우리가 원하던 근의 역할을 한다.

\begin{example}
$X^2+X+1$은 $\F_2[X]$에서 기약이다. 따라서
\[
  \F_2[X]/(X^2+X+1)
\]
은 네 원소를 갖는 체이다. $\alpha=X+(X^2+X+1)$라 하면
\[
  \alpha^2+\alpha+1=0,
  \qquad
  \alpha^2=\alpha+1
\]
이고, 모든 원소는
\[
  0,\ 1,\ \alpha,\ \alpha+1
\]
중 하나이다.
\end{example}

\section{같은 근, 서로 다른 표현}

근을 첨가한 체는 구체적인 수 안에서 찾을 수도 있고, 몫환으로 추상적으로 만들 수도 있다. 예를 들어
\[
  \Q[X]/(X^2-2)
\]
와 실수 안의 부분체 $\Q(\sqrt2)$는 서로 다른 방식으로 정의되지만 자연스럽게 동형이다. 대응은
\[
  f(X)+(X^2-2)\longmapsto f(\sqrt2)
\]
로 주어진다.

이 사실은 중요한 관점을 제공한다. 대수학에서 핵심은 원소가 ``실제로 무엇인가''보다 어떤 연산 관계를 만족하는가에 있다. 서로 다른 실현이 같은 대수적 구조를 가질 수 있다.

\begin{pitfall}
$K[X]/(f)$가 언제나 체인 것은 아니다. $f$가 기약일 때에만 $(f)$가 극대아이디얼이고 몫환이 체가 된다. 가약다항식으로 나누면 영인수가 생길 수 있다.
\end{pitfall}

\section{체가 얼마나 커졌는가}

$L$이 $K$를 포함하는 체이면 $L$은 자연스럽게 $K$-벡터공간이 된다. 따라서 체확장의 크기를
\[
  [L:K]=\dim_K L
\]
로 측정할 수 있다.

예를 들어
\[
  \Q(\sqrt2)
  =\{a+b\sqrt2:a,b\in\Q\}
\]
에서 $1,\sqrt2$는 $\Q$-기저이므로
\[
  [\Q(\sqrt2):\Q]=2.
\]
반면 $\R$은 $\Q$ 위에서 유한차원일 수 없다. 체확장의 차수는 단순히 원소의 개수를 세는 것이 아니라, 원래 체의 계수로 필요한 독립적인 방향의 수를 센다.

\section{체 안의 대칭}

체를 확장하면 새 원소를 움직이는 대칭이 나타난다. $L=\Q(\sqrt2)$에서 $\Q$의 원소를 고정하면서
\[
  \sqrt2\longmapsto-\sqrt2
\]
로 보내는 사상은 덧셈과 곱셈을 보존한다. 이 사상과 항등사상은 $L/\Q$의 자동동형군을 이룬다.

그러나 모든 체확장이 같은 수의 대칭을 가지는 것은 아니다. 예를 들어 $\Q(\sqrt[3]{2})$는 $X^3-2$의 한 실근을 포함하지만 두 비실근을 포함하지 않는다. 따라서 이 체 안에서는 $\sqrt[3]{2}$를 다른 근으로 보낼 수 없다. 모든 근을 담는 분해체로 이동해야 비로소 방정식의 전체 대칭이 드러난다.

\section{Volume 2의 지도}

이 권에서는 다음의 흐름을 따른다.

\begin{enumerate}
  \item 체의 표수와 가장 작은 부분체를 이해한다.
  \item 체확장을 벡터공간으로 보고 차수를 정의한다.
  \item 대수적 원소와 최소다항식을 연구한다.
  \item 몫체와 단순확장의 구조를 연결한다.
  \item 대수적 폐포와 분해체를 구성한다.
  \item 중근의 문제를 분리성과 비분리성으로 정리한다.
  \item 유한체의 구조와 프로베니우스 자동동형을 분석한다.
  \item 정상성과 분리성을 결합하여 갈루아 확장을 정의한다.
\end{enumerate}

\begin{checkpoint}
\begin{enumerate}[label=(\alph*)]
  \item $\R[X]/(X^2+1)$과 $\C$ 사이의 자연스러운 동형을 적어라.
  \item $\Q[X]/(X^2-2)$의 원소를 표준형으로 나타내라.
  \item $\Q(\sqrt[3]{2})$가 $X^3-2$의 모든 복소근을 포함하지 않는 이유를 설명하라.
\end{enumerate}
\end{checkpoint}

\section*{짧은 해설}

(a) $f(X)+(X^2+1)\mapsto f(i)$로 보낸다. (b) 모든 원소는 유일하게 $a+b\overline X$의 꼴이다. (c) 이 체는 실수체의 부분체인 반면 나머지 두 근은 비실수이기 때문이다.
